% !PRINTED 278
\weberpage{278}
%<<CONT
\[
(-1)^k x^k,
\]

and when the absolute value of $x$ is sufficiently large, the sign of this term will decide the sign of $L_k(x)$. If, then, we take for $\alpha$ a sufficiently large negative value, and for $\beta$ a sufficiently large positive value, then in (3) there occur only changes of sign for $x=\alpha$, and only permanences of sign for $x=\beta$. Hence $n$ roots will lie between $\alpha$ and $\beta$, and from this follows the theorem:

II.\ \gesperrt{The equation $L_n(x)=0$ has only real roots.}

But these two theorems are of only limited applicability so long as we cannot free ourselves from the assumption that in the chain of the $L_k$ no two consecutive terms are to vanish at the same time. From this restriction, however, we can free the theorem by the following simple consideration.

Suppose that for some value of $x$, $L_k$ vanishes but $L_{k-1}$ does not; we can then determine quantities $a_{k+1,i}$, which do not occur in $L_k$, in such a way that $L_{k+1}$ does not vanish for this value of $x$; for $L_{k+1}$ cannot vanish identically for all $a_{k+1,i}$, because it contains the term

\[
-a_{k+1,k}^2 L_{k-1}
\]

[\S.\,23, (12)].

Accordingly, if in the chain

\weq{5}{L_n, \; -L_{n-1}, \; L_{n-2} \cdots \pm L_1, \; \mp 1}

some consecutive terms vanish for some value of $x$, we can, by altering the $a_{i,k}$, derive another sequence

\weq{6}{L_n', \; -L_{n-1}', \; L_{n-2}' \cdots \pm L_1', \; \mp 1}

in which no two consecutive terms vanish for any value $x$ (of the interval $\alpha \ldots \beta$).

At the same time the $a_{i,k}'$, on which the $L'$ depend, can be so chosen that they \gesperrt{differ from the $a_{i,k}$ by less than an arbitrarily given quantity $\omega$}.

If now the numbers $\alpha, \beta$ are so chosen that in the sequence (5) no term vanishes for $x=\alpha$ or $x=\beta$, then we can take $\omega$ so small that corresponding terms of (5) and (6) have the same sign. But by our theorem the number of roots of $L_n'=0$ between $\alpha$ and $\beta$ is determined by the signs of the sequence (6).
